\section{Core Language}

\subsection{Codes, Capabilities, and Certificates}

\system separates the code a value currently inhabits from the rules that can act on that code.
In the running example, the four prepared data qubits begin with type $Q[\code{SurfaceD5},5,\State]$.
This type is a compact way to say three things the checker needs later: the qubit is encoded in the surface profile, the rule library treats it as distance five, and later operations may rely on state-level correctness rather than only final observable correctness.
The type is not a physical simulator state; it is the static handle that lets the compiler ask whether a requested operation is a direct capability, a switch target, or a resource-acquisition goal.

A code specification records the information needed to answer those questions.
The current prototype contains \code{SurfaceD5}, \code{Steane3}, \code{Tetra15}, and a negative-control \code{QLDPC12} profile.
\code{SurfaceD5} supports direct $\gate{H}$ and $\gate{S}$ in the small rule library; \code{Steane3} supports transversal Clifford gates; \code{Tetra15} supports transversal $\gate{T}$ and $\gate{CNOT}$; and \code{QLDPC12} requires long-range topology.
The last profile is deliberately included so that the checker has a concrete backend-incompatibility case to reject.

The capability context $\mathcal{K}$ collects the rule facts available to a compilation run:
\[
\begin{array}{rcl}
\mathcal{K} &::=& \Cap(C,g) \mid \Switch(C_1,C_2) \mid \Factory(R) \mid \Backend(B) \mid \cdots
\end{array}
\]
The interesting point is that these facts are not bare booleans.
Each fact carries a certificate level $\ell\in\{\Checked,\Certified,\Assumed\}$.
\Checked{} rules are discharged by local checkers or simple implementation invariants.
\Certified{} rules combine local checks with an explicit imported theorem from a cited protocol.
\Assumed{} rules are prototype glue and are rejected by the strict theorem policy.
This is the mechanism that prevents the paper from confusing a useful compiler experiment with a new physical proof.

\subsection{Source and Target Syntax}

The source language is intentionally ordinary because users should describe logical intent rather than acquisition mechanics.
It includes preparation, logical gates, error-correction rounds, logical measurement, sequencing, and classical branching:
\[
\begin{array}{rcl}
e &::=& \mathsf{prepare}\ q\ s
  \mid \mathsf{gate}\ g\ \vec q
  \mid \mathsf{ec}\ \vec q\ d
  \mid \mathsf{measure}\ b\ O\ \vec q \\
  && \mid\ \mathsf{if}\ b\ \mathsf{then}\ e_1\ \mathsf{else}\ e_2
  \mid e_1;e_2.
\end{array}
\]
In the case study, this language is represented by \texttt{LogicalOp.prepare}, \texttt{LogicalOp.apply}, \texttt{LogicalOp.ec}, \texttt{LogicalOp.measure}, and \texttt{LogicalOp.branch}.
The source contains $\gate{T}$ and $\gate{CNOT}$ as logical gates even though the initial code does not directly support them.
That is the point of elaboration: unsupported source gates become acquisition obligations rather than immediate type errors.

The target language is smaller and more explicit.
It exposes the atoms that a checker can audit:
\[
\begin{array}{rcl}
a &::=& \mathsf{alloc}\ q:C
  \mid \mathsf{native}\ g\ \vec q
  \mid \mathsf{transv}\ g\ \vec q
  \mid \mathsf{switch}\ q:C_1\to C_2 \\
  && \mid\ \mathsf{produce}\ r:R
  \mid \mathsf{consume}\ r\ \mathsf{for}\ g\ \vec q
  \mid \mathsf{measureL}\ b\ O\ \vec q \\
  && \mid\ \mathsf{decode}\ s
  \mid \mathsf{frame}\ f
  \mid \mathsf{postselect}\ p
  \mid \mathsf{reset\_clean}\ q, \\
t &::=& a \mid t_1;t_2 \mid \mathsf{if}\ b\ \mathsf{then}\ t_1\ \mathsf{else}\ t_2.
\end{array}
\]
The target has no unrestricted $\mathsf{apply}\ U$ primitive.
That omission is a design constraint, not an aesthetic preference.
If arbitrary unitaries could appear as primitive atoms, a compiler could hide an entire code-switching or factory protocol behind one step, and the checker would lose the ability to inspect capability, resource, effect, and certificate evidence.

\subsection{Types, Resources, and Effects}

The encoded quantum type $Q[C,d,\mu]$ is the thread that connects the paper calculus to the implementation.
The code index $C$ changes when the planner inserts a switch, stays fixed for direct and resource-injection gates, and must agree across branches before the program continues.
The mode $\mu ::= \State \mid \Observable$ distinguishes state-level values from values whose guarantee is only about final measurement distributions.
The implementation already carries this field in \texttt{QType}; the calculus uses it to state the intended discipline that observable-level outputs should not be fed to rules requiring state-level inputs.

Probabilistic resources live in a separate context because they are physical artifacts, not reusable descriptions:
\[
R ::= \res{TState}[\epsilon] \mid \res{CCZState}[\epsilon] \mid \res{BellPair}[\epsilon] \mid \res{AuxState}[C,\epsilon].
\]
In the implementation, a resource acquisition step places a fresh identifier such as \texttt{res\_t\_0} in the sidecar, and the corresponding consume step names the same identifier.
The checker maintains produced and consumed sets while flattening the plan.
This small piece of bookkeeping is the operational version of the linear resource context $\Delta$.

Effects summarize what a plan promises conservatively:
\[
E = \langle \epsilon, f, a, s, \tau, \rho, n_{\mathsf{sw}}, n_{\mathsf{fac}}, \ldots\rangle.
\]
Here $\epsilon$ is a logical-error upper bound, $f$ is a failure upper bound, $a$ is an acceptance lower bound, $s$ is peak space, $\tau$ is cycle count, $\rho$ is qubit-round count, and the remaining fields count switches, factories, measurements, resets, multi-qubit gates, decoder latency, certificates, assumptions, and rule identifiers.
Sequential composition adds error, cycles, qubit-rounds, and counts, multiplies acceptance lower bounds, and takes a maximum for peak space.
Branch composition takes conservative branchwise bounds.
The effect algebra is deliberately simple because it must be checked on every generated sidecar; its value is that the numbers in Section~\ref{sec:case-study} are produced by the same composition discipline used by the checker.
